\documentclass[11pt]{article}

\usepackage[margin=1in]{geometry}
\usepackage{amsmath, amssymb}
\usepackage[T1]{fontenc}
\usepackage{lmodern}
\usepackage{microtype}
\usepackage{graphicx}
\usepackage{booktabs}
\usepackage{array}
\usepackage{float}
\usepackage{xcolor}
\usepackage{fancyhdr}
\usepackage[font=small,labelfont=bf,skip=5pt]{caption}
\usepackage{listings}
\lstdefinestyle{pycode}{
  language=Python, basicstyle=\ttfamily\scriptsize,
  keywordstyle=\color{psink}\bfseries, commentstyle=\color{psaccent}\itshape,
  stringstyle=\color{psink}, numbers=left, numberstyle=\tiny\color{psrule},
  numbersep=7pt, breaklines=true, breakatwhitespace=true, showstringspaces=false,
  frame=none, xleftmargin=20pt, aboveskip=6pt, belowskip=6pt, columns=fullflexible,
  keepspaces=true, upquote=true,
}
\lstset{style=pycode, literate={²}{{\textsuperscript{2}}}1}

\definecolor{psink}{RGB}{23,54,93}
\definecolor{psaccent}{RGB}{192,80,77}
\definecolor{psrule}{RGB}{200,209,222}

\usepackage[colorlinks=true,allcolors=psink]{hyperref}


\emergencystretch=3em

\setlength{\textfloatsep}{10pt plus 2pt minus 2pt}
\setlength{\intextsep}{10pt plus 2pt minus 2pt}
\setlength{\abovecaptionskip}{4pt}
\setlength{\belowcaptionskip}{0pt}

\makeatletter
\renewcommand\section{\@startsection{section}{1}{0pt}%
  {-2.6ex \@plus -1ex \@minus -.2ex}{1.3ex \@plus .2ex}%
  {\normalfont\scshape\large\color{psink}}}   % lmr has no bold small caps
\renewcommand\subsection{\@startsection{subsection}{2}{0pt}%
  {-2.2ex \@plus -1ex \@minus -.2ex}{1.0ex \@plus .2ex}%
  {\normalfont\bfseries\normalsize\color{psink}}}
\renewcommand\subsubsection{\@startsection{subsubsection}{3}{0pt}%
  {-1.9ex \@plus -.8ex \@minus -.2ex}{0.8ex \@plus .2ex}%
  {\normalfont\bfseries\itshape\normalsize\color{psink}}}
\renewcommand\subparagraph{\@startsection{subparagraph}{5}{0pt}%
  {-1.6ex \@plus -.6ex \@minus -.2ex}{0.5ex \@plus .1ex}%
  {\normalfont\itshape\normalsize\color{psink}}}
\renewcommand\paragraph{\@startsection{paragraph}{4}{0pt}%
  {1.7ex \@plus 1ex \@minus .2ex}{-0.9em}%
  {\normalfont\bfseries\normalsize\color{psink}}}
\makeatother

\newenvironment{tight}
  {\begin{itemize}\setlength{\itemsep}{3pt}\setlength{\parskip}{0pt}%
   \setlength{\topsep}{3pt}\setlength{\leftmargini}{1.4em}}
  {\end{itemize}}

\setcounter{secnumdepth}{0}   % headings carry the assignment's own labels
\setlength{\headheight}{14pt}
\pagestyle{fancy}
\fancyhf{}
\fancyhead[L]{\footnotesize\scshape\color{psink} MGT 595 \textemdash\ Problem Set 4}
\fancyhead[R]{\footnotesize\color{psink}\thepage}
\renewcommand{\headrulewidth}{0.4pt}
\renewcommand{\headrule}{\hbox to\headwidth{%
  \color{psrule}\leaders\hrule height \headrulewidth\hfill}}

\begin{document}

\thispagestyle{empty}
\begin{flushleft}
{\bfseries\LARGE\color{psink} Quantitative Investing}\\[3pt]
{\scshape\Large\color{psink} Problem Set 4}\\[10pt]
{\color{psrule}\rule{\textwidth}{1.2pt}}\\[5pt]
{\small Ashley Wen\,(xw564)\quad\textbar\quad Emily Wu\,(esw63)%
\quad\textbar\quad Grace Gu\,(xg325)\hfill \today}
\end{flushleft}

\vspace{1.4em}

{\color{psink}\textbf{Contents.}} Parts I--III test the CAPM in the
time-series framework with the Gibbons, Ross and Shanken (GRS) test, on the
30 value-weight industry portfolios, the 10 past return portfolios and the 25
size and BE/ME portfolios. Part III also examines Roll's critique using
in-sample and out-of-sample tangency portfolios as the market proxy. The
capstone adds the profitability factor RMW to the market and tests the
augmented model on the 25 size and BE/ME portfolios. The capstone Python
code is reproduced in its section, in the order of the original notebook.

\vspace{0.8em}


% =====================================================================
\section{Part I: 30 Value-Weight Industry Portfolios}
% =====================================================================

\paragraph{Data and conventions.} Returns are monthly percentages. Portfolio
returns are converted to excess returns by subtracting RF; RM-RF is already
an excess return.

\subsection{Question (a)}

\begin{table}[H]
\centering
\caption{Descriptive statistics, 30 value-weight industry portfolios.}
\footnotesize
\setlength{\tabcolsep}{4pt}
\begin{tabular}{lrrrclrrr}
\toprule
Portfolio & Mean & SD & Sharpe && Portfolio & Mean & SD & Sharpe \\
\midrule
Food & 0.940 & 4.685 & 0.143 && Carry & 1.153 & 7.577 & 0.116 \\
Beer & 1.135 & 7.038 & 0.123 && Mines & 0.979 & 7.356 & 0.096 \\
Smoke & 1.167 & 5.841 & 0.154 && Coal & 1.075 & 10.998 & 0.073 \\
Games & 1.117 & 8.945 & 0.095 && Oil & 1.034 & 6.423 & 0.119 \\
Books & 0.889 & 7.165 & 0.086 && Util & 0.887 & 5.435 & 0.114 \\
Hshld & 0.891 & 5.737 & 0.108 && Telcm & 0.812 & 4.661 & 0.116 \\
Clths & 0.871 & 6.225 & 0.096 && Servs & 1.237 & 8.693 & 0.111 \\
Hlth & 1.068 & 5.501 & 0.145 && BusEq & 1.237 & 6.745 & 0.143 \\
Chems & 1.020 & 6.307 & 0.119 && Paper & 0.964 & 5.861 & 0.118 \\
Txtls & 0.927 & 7.936 & 0.083 && Trans & 0.929 & 7.068 & 0.093 \\
Cnstr & 1.030 & 6.973 & 0.109 && Whlsl & 0.869 & 7.247 & 0.082 \\
Steel & 0.997 & 8.637 & 0.084 && Rtail & 1.061 & 5.951 & 0.133 \\
FabPr & 1.149 & 7.268 & 0.121 && Meals & 1.065 & 6.433 & 0.123 \\
ElcEq & 1.202 & 7.700 & 0.121 && Fin & 1.010 & 6.724 & 0.110 \\
Autos & 1.209 & 8.578 & 0.109 && Other & 0.821 & 6.582 & 0.084 \\
\bottomrule
\end{tabular}
\label{tab:p1a}
\vspace{2pt}\begin{minipage}{0.93\textwidth}\footnotesize
Notes: Monthly data, July 1926--June 2026. Mean and standard deviation are
based on raw portfolio returns, in percent per month. Sharpe ratios are
computed using excess returns, $(R_p-R_f)$, divided by their sample standard
deviation.
\end{minipage}
\end{table}

There is no common ordering across industries, but two patterns are visible.
Average returns are only weakly positively related to market beta
(correlation about 0.30), while Sharpe ratios are negatively related to beta
(correlation about $-0.52$). Only five industries have monthly Sharpe ratios
above the market's 0.131: Smoke, Health, Food, Business Equipment, and
Retail.

\subsection{Question (b)}

For each industry portfolio, we estimated
\[
R_{it}-R_{ft} = \alpha_i+\beta_i(R_{Mt}-R_{ft})+\varepsilon_{it}.
\]
The sample is July 1926 through June 2026, with $T=1200$, $N=30$, and $K=1$.

\begin{table}[H]
\centering
\caption{GRS test, 30 industry portfolios on RM-RF.}
\small
\begin{tabular}{lr}
\toprule
Statistic & Result \\
\midrule
GRS F-statistic & 1.7270 \\
Degrees of freedom & $F(30, 1169)$ \\
$p$-value & 0.0091 \\
Individually significant alphas at 5\% & 4 of 30 \\
\bottomrule
\end{tabular}
\label{tab:p1b_grs}
\end{table}

Because the $p$-value is below 1\%, we reject the joint null that all 30
industry alphas equal zero. The one-factor CAPM does not fully price the
industry portfolios.

\begin{table}[H]
\centering
\caption{CAPM time-series regressions, 30 industry portfolios.}
\footnotesize
\setlength{\tabcolsep}{4pt}
\begin{tabular}{lrrrrclrrrr}
\toprule
Portfolio & Alpha & Beta & $t(\alpha)$ & $p(\alpha)$ && Portfolio & Alpha & Beta & $t(\alpha)$ & $p(\alpha)$ \\
\midrule
Food & 0.167 & 0.723 & 2.14 & 0.0328 && Carry & 0.061 & 1.182 & 0.49 & 0.6262 \\
Beer & 0.232 & 0.911 & 1.55 & 0.1211 && Mines & 0.075 & 0.911 & 0.47 & 0.6418 \\
Smoke & 0.464 & 0.624 & 3.31 & 0.0009 && Coal & $-$0.078 & 1.269 & $-$0.31 & 0.7594 \\
Games & $-$0.120 & 1.390 & $-$0.81 & 0.4175 && Oil & 0.152 & 0.879 & 1.18 & 0.2378 \\
Books & $-$0.152 & 1.109 & $-$1.28 & 0.2008 && Util & 0.090 & 0.759 & 0.85 & 0.3979 \\
Hshld & 0.013 & 0.875 & 0.13 & 0.8977 && Telcm & 0.079 & 0.666 & 0.90 & 0.3708 \\
Clths & 0.016 & 0.842 & 0.13 & 0.8996 && Servs & 0.352 & 0.883 & 1.65 & 0.0991 \\
Hlth & 0.223 & 0.827 & 2.30 & 0.0215 && BusEq & 0.215 & 1.082 & 2.07 & 0.0391 \\
Chems & 0.025 & 1.042 & 0.28 & 0.7762 && Paper & 0.036 & 0.947 & 0.40 & 0.6858 \\
Txtls & $-$0.146 & 1.156 & $-$0.99 & 0.3205 && Trans & $-$0.139 & 1.148 & $-$1.33 & 0.1850 \\
Cnstr & $-$0.065 & 1.186 & $-$0.74 & 0.4597 && Whlsl & $-$0.164 & 1.097 & $-$1.30 & 0.1949 \\
Steel & $-$0.227 & 1.372 & $-$1.67 & 0.0958 && Rtail & 0.117 & 0.970 & 1.33 & 0.1833 \\
FabPr & 0.014 & 1.245 & 0.15 & 0.8780 && Meals & 0.140 & 0.943 & 1.18 & 0.2375 \\
ElcEq & 0.029 & 1.299 & 0.29 & 0.7697 && Fin & $-$0.064 & 1.156 & $-$0.79 & 0.4311 \\
Autos & 0.038 & 1.296 & 0.25 & 0.7992 && Other & $-$0.167 & 1.033 & $-$1.57 & 0.1169 \\
\bottomrule
\end{tabular}
\label{tab:p1b_reg}
\vspace{2pt}\begin{minipage}{0.93\textwidth}\footnotesize
Notes: Alphas are monthly percentage returns. Alpha $t$-statistics use the OLS
standard errors from the corresponding time-series CAPM regression.
\end{minipage}
\end{table}

\paragraph{AI integration: code generation --- GRS F-statistic.}

\textit{Prompt 1.} ``State the GRS F-statistic formula precisely and explain
each term.''

The LLM correctly identified the main structure of the GRS statistic,
including the joint alpha term, the factor mean-variance adjustment, and the
$F(N,T-N-K)$ distribution. However, the LLM used the original GRS
normalization, in which the covariance matrices are normalized differently
from the convention specified in this assignment. The assignment explicitly
requires the residual covariance matrix to use denominator $T-K-1$, the
factor covariance matrix to use denominator $T-1$, and the GRS statistic to
use the scaling $(T-N-K)/N$. We therefore followed the assignment's stated
convention exactly. The difference is a normalization choice and does not
change the qualitative conclusion in this sample. For this problem set, we
used
\[
F_{GRS} = \frac{T-N-K}{N}\,
\frac{\hat{\alpha}'\hat{\Sigma}_{\varepsilon}^{-1}\hat{\alpha}}
{1+\bar f'\hat{\Omega}_f^{-1}\bar f}.
\]

\textit{Prompt 2.} ``Using the GRS formula above, write Python code
implementing the GRS test for 30 industry portfolios in a one-factor CAPM.
Use residual covariance with denominator $T-K-1$, factor covariance with
denominator $T-1$, and report the GRS F-statistic and $p$-value.''

The LLM correctly used residual covariance rather than raw-return covariance
and correctly recognized the required F-distribution degrees of freedom.
However, it retained the original-GRS normalization by adding an extra
finite-sample scaling adjustment. Because this differs from the convention
explicitly specified in the assignment, we used the assignment formula
exactly. Using the assignment convention, the 30-industry test gives a GRS
F-statistic of \textbf{1.7270}, a $p$-value of \textbf{0.0091}, and degrees
of freedom \textbf{F(30, 1169)}.

\textit{Validation checks.} We performed the three validation checks required
in the assignment:
\begin{enumerate}
\setlength{\itemsep}{2pt}\setlength{\parskip}{0pt}
\item \textbf{Zero-alpha check:} Setting all alphas equal to zero produced a
GRS statistic of exactly zero.
\item \textbf{Reorder invariance:} Permuting the order of the 30 test
portfolios did not change the GRS statistic.
\item \textbf{Scale invariance:} Multiplying all returns by 100 did not
change the GRS statistic.
\end{enumerate}
All three checks passed. As an additional cross-check, we verified the
one-factor GRS statistic using the equivalent Sharpe-ratio formulation,
which reproduced the same result to numerical precision.

\subsection{Question (c)}

The GRS null is
\[
H_0:\alpha_1=\alpha_2=\cdots=\alpha_{30}=0.
\]
Under the CAPM, expected excess returns should be fully explained by
exposure to the market factor, so every time-series intercept should be
zero. Equivalently, the market proxy should be mean-variance efficient
relative to the test assets.

Intuitively, GRS asks whether the alphas are jointly too large to be
attributed to sampling noise. It combines the pricing errors while
accounting for their residual variances and cross-sectional residual
correlations. Large, precisely estimated pricing errors that are not
redundant with one another contribute more strongly to rejection.

The Sharpe-ratio interpretation is equivalent: if the market proxy is
mean-variance efficient, adding the test assets should not significantly
improve the maximum attainable Sharpe ratio. Rejecting GRS provides evidence
that the market proxy lies inside, rather than on, the mean-variance
frontier spanned by the factor and test assets.

The time-series regressions also implicitly impose the CAPM beta risk
premium. Since the factor is RM-RF, its sample mean is the estimated market
risk premium. The CAPM predicts
\[
E[R_i-R_f]=\beta_iE[R_M-R_f],
\]
so the intercept $\alpha_i$ measures the portfolio's deviation from this
pricing relation.

\paragraph{AI integration: null hypothesis and the Sharpe ratio
interpretation.}

\textit{Prompt 1.} ``State the null hypothesis of the GRS test precisely.
What does rejecting it imply about the CAPM?''

The LLM correctly stated the joint null hypothesis
$H_0:\alpha_1=\alpha_2=\cdots=\alpha_N=0$. It also correctly explained that
rejecting the null means the CAPM fails to price the test portfolios
jointly. The LLM correctly noted that rejection does not imply that every
individual alpha must be statistically significant. However, it did not
explicitly state the equivalent mean-variance interpretation required in the
assignment: under the null, the market proxy is mean-variance efficient with
respect to the test assets. The LLM also did not explicitly explain the
final point required in the assignment: in the time-series CAPM, the beta
risk premium is implicitly set equal to the sample mean of RM-RF.
Therefore, the CAPM pricing relation is
$E[R_i-R_f]=\beta_iE[R_M-R_f]$, and the intercept $\alpha_i$ measures the
deviation from that relation.

\textit{Prompt 2.} ``The GRS test statistic can be interpreted in terms of
Sharpe ratios. Specifically, it is related to whether the test assets can
improve the Sharpe ratio of the factor portfolio used as the market proxy.
Can you explain this interpretation?''

The LLM correctly explained that the GRS test is related to the improvement
in the maximum squared Sharpe ratio that becomes possible when the test
assets are added to the factor portfolio. It correctly connected
$\alpha'\Sigma_{\varepsilon}^{-1}\alpha$ to the mean-variance value of the
pricing errors and explained that rejecting GRS means there is statistically
significant evidence that adding the test assets can improve the attainable
Sharpe ratio. It also correctly connected rejection to the market proxy
lying inside, rather than on, the mean-variance frontier spanned by the
factor and test assets. The conceptual interpretation was correct. However,
the response again used the original-GRS finite-sample normalization rather
than the convention specified in this assignment, so we retained the
assignment normalization in the implementation and reported results.

\subsection{Question (d)}

Most industry pricing errors are economically small: the average absolute
alpha is about 0.13\% per month, and 26 of the 30 individual alphas are not
significant at the 5\% level. However, the alphas are not randomly
distributed. The correlation between beta and alpha is approximately $-0.67$.

Lower-beta industries such as Smoke ($\beta=0.624$, $\alpha=0.464\%$), Food
($\beta=0.723$, $\alpha=0.167\%$), and Health ($\beta=0.827$,
$\alpha=0.223\%$) have positive alphas, while high-beta industries such as
Steel ($\beta=1.372$, $\alpha=-0.227\%$) and Games ($\beta=1.390$,
$\alpha=-0.120\%$) have negative alphas.

Business Equipment is an exception to the simple low-beta pattern: it has
$\beta=1.082$ but still has a significantly positive alpha of 0.215\% per
month. Overall, however, the negative beta-alpha relationship is consistent
with a flatter security market line than the CAPM predicts. One possible
explanation is that market beta alone does not capture all dimensions of
risk or expected returns across industries. An imperfect market proxy could
also contribute to the pricing errors. The data here show that the
one-factor CAPM misses the pattern, although they do not by themselves
distinguish among these possible explanations.


% =====================================================================
\section{Part II: 10 Past Return Portfolios}
% =====================================================================

\subsection{Question (e)}

The momentum sample is January 1927 through June 2026, with $T=1194$.

\paragraph{Descriptive statistics.}

\begin{table}[H]
\centering
\caption{Descriptive statistics, 10 past return portfolios.}
\small
\begin{tabular}{lrrr}
\toprule
Portfolio & Mean raw return (\%/mo) & Raw return SD (\%/mo) & Monthly Sharpe (excess) \\
\midrule
Loser & 0.372 & 9.769 & 0.010 \\
2 & 0.732 & 8.037 & 0.057 \\
3 & 0.793 & 6.879 & 0.076 \\
4 & 0.908 & 6.296 & 0.101 \\
5 & 0.937 & 5.841 & 0.114 \\
6 & 0.984 & 5.721 & 0.125 \\
7 & 1.026 & 5.394 & 0.140 \\
8 & 1.123 & 5.278 & 0.161 \\
9 & 1.171 & 5.529 & 0.163 \\
Winner & 1.534 & 6.498 & 0.194 \\
\bottomrule
\end{tabular}
\label{tab:p2e_stats}
\vspace{2pt}\begin{minipage}{0.93\textwidth}\footnotesize
Notes: Monthly data, January 1927--June 2026. Mean and standard deviation
are based on raw portfolio returns. Sharpe ratios are computed using excess
returns, $(R_p-R_f)$, divided by their sample standard deviation.
\end{minipage}
\end{table}

Unlike the industry portfolios, the past-return portfolios show a clear
monotonic pattern. Average returns and Sharpe ratios rise strongly from past
losers to past winners. The monthly Sharpe ratio rises from about 0.01 for
Loser to about 0.19 for Winner.

\paragraph{CAPM and GRS results.}

\begin{table}[H]
\centering
\caption{GRS test, 10 past return portfolios on RM-RF.}
\small
\begin{tabular}{lr}
\toprule
Statistic & Result \\
\midrule
GRS F-statistic & 6.5930 \\
Degrees of freedom & $F(10, 1183)$ \\
$p$-value & $5.34\times10^{-10}$ \\
Individually significant alphas at 5\% & 6 of 10 \\
\bottomrule
\end{tabular}
\label{tab:p2e_grs}
\end{table}

The GRS test strongly rejects the joint null that all ten momentum alphas
equal zero. The rejection is substantially stronger than for the industry
portfolios.

\begin{table}[H]
\centering
\caption{CAPM time-series regressions, 10 past return portfolios.}
\small
\begin{tabular}{lrrrr}
\toprule
Portfolio & Alpha (\%/mo) & Beta & Alpha $t$ & Alpha $p$ \\
\midrule
Loser & $-$0.977 & 1.559 & $-$6.43 & 0.0000 \\
2 & $-$0.462 & 1.335 & $-$4.17 & 0.0000 \\
3 & $-$0.286 & 1.169 & $-$3.31 & 0.0010 \\
4 & $-$0.120 & 1.094 & $-$1.68 & 0.0924 \\
5 & $-$0.048 & 1.032 & $-$0.80 & 0.4236 \\
6 & 0.004 & 1.026 & 0.07 & 0.9440 \\
7 & 0.090 & 0.961 & 1.75 & 0.0809 \\
8 & 0.207 & 0.933 & 3.87 & 0.0001 \\
9 & 0.239 & 0.957 & 3.74 & 0.0002 \\
Winner & 0.552 & 1.028 & 5.36 & 0.0000 \\
\bottomrule
\end{tabular}
\label{tab:p2e_reg}
\vspace{2pt}\begin{minipage}{0.93\textwidth}\footnotesize
Notes: Alphas are monthly percentage returns. Alpha $t$-statistics use the
OLS standard errors from the corresponding time-series CAPM regression.
\end{minipage}
\end{table}

The CAPM has particular difficulty explaining the winner-loser spread.
Across the ten portfolios, the correlation between market beta and average
excess return is approximately $-0.82$, and the correlation between beta and
alpha is approximately $-0.91$.

The Loser portfolio has the highest beta, $\beta=1.559$, but the lowest
average return and a large negative alpha of $-0.977\%$ per month. The
Winner portfolio has a much lower beta, $\beta=1.028$, but a much higher
average return and a positive alpha of $0.552\%$ per month.

Thus, market beta points in the wrong direction for explaining the momentum
return spread. Past losers tend to have negative alphas and past winners
positive alphas. The winner and loser alphas have opposite signs, although
the magnitudes are not symmetric because the loser-side pricing error is
larger.

One interpretation is that momentum captures exposure to an omitted priced
risk factor that is not included in the CAPM. Another is behavioral
underreaction or slow information diffusion, which can create predictable
continuation in returns. The CAPM itself cannot distinguish between these
explanations.

\paragraph{AI integration: momentum and the CAPM.}

\textit{Prompt.} ``The CAPM typically rejects more strongly for portfolios
sorted on past returns than for industry portfolios. Why is momentum
particularly difficult for the CAPM to price? Is momentum evidence against
the CAPM, against efficient markets, or both?''

The LLM correctly explained that the CAPM has no direct mechanism for
momentum: market beta cannot explain why past winners subsequently earn much
higher returns than past losers. It also correctly distinguished two broad
interpretations:
\begin{itemize}
\setlength{\itemsep}{2pt}\setlength{\parskip}{0pt}
\item \textbf{Risk-based explanation:} momentum may reflect exposure to an
omitted priced factor, so the CAPM is incomplete while market efficiency
could still hold.
\item \textbf{Behavioral explanation:} underreaction, slow information
diffusion, or other investor behavior may create persistent mispricing.
\end{itemize}
The LLM therefore correctly concluded that momentum is strong evidence
against the CAPM, but is not automatically evidence against market
efficiency because of the joint-hypothesis problem.

However, part of its numerical intuition was too generic. It suggested that
winner and loser portfolios often have similar market betas. This does not
match our sample. In our data, Loser beta $=$ \textbf{1.559}, Winner beta
$=$ \textbf{1.028}, the correlation between beta and average excess return
across the ten portfolios is approximately \textbf{$-$0.82}, and the
correlation between beta and alpha across the ten portfolios is
approximately \textbf{$-$0.91}. Thus, higher market beta is actually
associated with lower average returns and more negative alphas across the
momentum portfolios.

The alpha pattern does match the LLM's general prediction: Loser alpha $=$
\textbf{$-$0.977\% per month} and Winner alpha $=$ \textbf{+0.552\% per
month}. The winner and loser alphas therefore have opposite signs, although
their magnitudes are not symmetric because the loser-side pricing error is
larger. The LLM correctly distinguished the risk-based and behavioral
explanations, but its discussion was mostly conceptual and did not provide
much empirical evidence that would distinguish between the two.

\paragraph{Bonus: time-series vs.\ cross-sectional momentum.}

\textit{Prompt.} ``Are time-series momentum (TSMOM) and cross-sectional
momentum (UMD) the same thing? Explain the difference.''

The LLM correctly explained that the two strategies are related but not
identical. Cross-sectional momentum ranks assets relative to one another and
goes long relative winners and short relative losers. UMD is an example of
this type of momentum strategy. Time-series momentum instead uses the sign
of each asset's own past return: an asset with positive past performance is
typically held long, while an asset with negative past performance is held
short. Therefore, cross-sectional momentum uses relative performance across
assets, while time-series momentum uses each asset's own return history. The
two strategies can even take opposite positions in the same asset.


% =====================================================================
\section{Part III: 25 Size and BE/ME Portfolios}
% =====================================================================

\paragraph{Sample.} 1926-07 to 2026-06, $T = 1200$, $N = 25$. Returns are in excess of the one-month T-bill rate, in percent per month (Parts I and II report raw means; Sharpe ratios use excess returns throughout). Rows sort on size
(S = small, B = big); columns sort on BE/ME (L = low, H = high).

\subsection{Question (f), repeating Part I(a)}

\begin{table}[H]
\centering
\caption{Summary statistics, 25 size and BE/ME portfolios, percent per month.}
\small
\begin{tabular}{lrrrrrcrrrrr}
\toprule
& \multicolumn{5}{c}{Mean excess return} && \multicolumn{5}{c}{Standard deviation} \\
\cmidrule{2-6}\cmidrule{8-12}
& L & 2 & 3 & 4 & H && L & 2 & 3 & 4 & H \\
\midrule
S & 0.58 & 0.68 & 1.00 & 1.14 & 1.35 && 11.88 & 9.62 & 8.82 & 8.25 & 9.13 \\
2 & 0.66 & 0.95 & 0.98 & 1.04 & 1.24 &&  7.96 & 7.46 & 7.21 & 7.41 & 8.64 \\
3 & 0.73 & 0.91 & 0.92 & 1.01 & 1.11 &&  7.34 & 6.47 & 6.45 & 6.91 & 8.34 \\
4 & 0.73 & 0.79 & 0.86 & 0.97 & 1.03 &&  6.17 & 6.08 & 6.28 & 6.80 & 8.52 \\
B & 0.69 & 0.66 & 0.73 & 0.67 & 1.01 &&  5.34 & 5.24 & 5.56 & 6.50 & 8.48 \\
\midrule
& \multicolumn{5}{c}{Sharpe ratio} && \multicolumn{5}{c}{} \\
\cmidrule{2-6}
S & 0.049 & 0.071 & 0.113 & 0.138 & 0.148 &&&&&& \\
2 & 0.083 & 0.127 & 0.137 & 0.141 & 0.144 &&&&&& \\
3 & 0.100 & 0.141 & 0.142 & 0.145 & 0.133 &&&&&& \\
4 & 0.118 & 0.129 & 0.136 & 0.143 & 0.121 &&&&&& \\
B & 0.130 & 0.125 & 0.131 & 0.103 & 0.119 &&&&&& \\
\bottomrule
\end{tabular}
\label{tab:p3f_stats}
\vspace{2pt}\begin{minipage}{0.93\textwidth}\footnotesize
Monthly Sharpe ratio $= \bar R^e_i/\hat\sigma_i$. Market proxy (RM$-$RF):
mean $0.70$, standard deviation $5.31$, Sharpe ratio $0.131$.
\end{minipage}
\end{table}

\paragraph{Discussion.}
\begin{itemize}
\setlength{\itemsep}{4pt}\setlength{\parskip}{0pt}

\item \textbf{Is there a discernible pattern in the average returns or
Sharpe ratios?} The dominant pattern is a value effect. Average returns and
Sharpe ratios rise with BE/ME in every size quintile, and the spread is
steepest among small stocks. The size effect is conditional: small value
beats big value, but small growth earns less than big growth and is the
worst of the 25 portfolios, with the lowest mean and the highest
volatility. Unlike the industries in Part I, which span a similar range of Sharpe ratios with no ordering, the dispersion here lines up with the characteristic we sorted on, which is exactly what a single beta will struggle to match.

\end{itemize}

\subsection{Question (f), repeating Part I(b)}

For each portfolio we estimate
\[
R_{it} - R_{ft} = \alpha_i + \beta_{iM}\,(R_{Mt} - R_{ft}) + e_{it},
\]
and test $H_0:\alpha_1 = \dots = \alpha_{25} = 0$ with the GRS statistic,
distributed $F(N,\,T-N-K) = F(25,\,1174)$ under the null.

\begin{table}[H]
\centering
\caption{CAPM time-series regressions, 25 size and BE/ME portfolios.
Intercepts in percent per month.}
\small
\begin{tabular}{lrrrrrcrrrrr}
\toprule
& \multicolumn{5}{c}{$\hat\alpha_i$} && \multicolumn{5}{c}{$t(\hat\alpha_i)$} \\
\cmidrule{2-6}\cmidrule{8-12}
& L & 2 & 3 & 4 & H && L & 2 & 3 & 4 & H \\
\midrule
S & $-0.53$ & $-0.29$ & 0.06 & 0.26 & 0.41 && $-2.20$ & $-1.59$ & 0.38 & 1.88 & 2.48 \\
2 & $-0.22$ & 0.09 & 0.15 & 0.20 & 0.29 && $-1.80$ & 0.88 & 1.53 & 1.88 & 2.12 \\
3 & $-0.13$ & 0.13 & 0.14 & 0.19 & 0.17 && $-1.43$ & 1.77 & 1.88 & 2.17 & 1.37 \\
4 & $-0.03$ & 0.03 & 0.09 & 0.17 & 0.07 && $-0.41$ & 0.50 & 1.36 & 1.96 & 0.57 \\
B & 0.02 & 0.00 & 0.06 & $-0.09$ & 0.10 && 0.51 & 0.01 & 0.90 & $-1.05$ & 0.72 \\
\midrule
& \multicolumn{5}{c}{$\hat\beta_{iM}$} && \multicolumn{5}{c}{} \\
\cmidrule{2-6}
S & 1.60 & 1.39 & 1.36 & 1.26 & 1.36 &&&&&& \\
2 & 1.27 & 1.23 & 1.20 & 1.21 & 1.37 &&&&&& \\
3 & 1.25 & 1.13 & 1.12 & 1.17 & 1.35 &&&&&& \\
4 & 1.09 & 1.09 & 1.10 & 1.16 & 1.39 &&&&&& \\
B & 0.96 & 0.94 & 0.96 & 1.09 & 1.30 &&&&&& \\
\bottomrule
\end{tabular}
\label{tab:p3f_capm}
\vspace{2pt}\begin{minipage}{0.93\textwidth}\footnotesize
GRS $= 3.37$, $p < 0.001$, $F(25, 1174)$. $\hat\Sigma$ uses denominator
$T-K-1$; the factor variance uses $T-1$.
\end{minipage}
\end{table}

\begin{itemize}
\setlength{\itemsep}{4pt}\setlength{\parskip}{0pt}

\item \textbf{What do the GRS F-statistic and its $p$-value show?} The GRS test rejects ($F = 3.37$, $p < 0.001$), so the market proxy is not mean--variance efficient with respect to the 25 size and BE/ME
portfolios.

\end{itemize}

\subsection{Question (f), repeating Part I(d)}

\paragraph{Discussion.}
\begin{itemize}
\setlength{\itemsep}{4pt}\setlength{\parskip}{0pt}

\item \textbf{What are the sign and magnitude of the intercepts?} The
intercepts are negative for growth and positive for value. The spread is widest among small stocks, from about $-0.5\%$ per
month for small growth to $+0.4\%$ for small value. Among the largest
stocks the intercepts are close to zero.

\item \textbf{Does the CAPM have difficulty pricing any portfolios in
particular?} The CAPM has most difficulty with small growth and small value.
Small growth is overpriced and small value underpriced relative to
their betas, while big stocks are priced reasonably well. Only five
intercepts are significant on their own, but the residuals are strongly
correlated across portfolios, so the GRS test reads the alphas as one
systematic pattern rather than 25 small deviations.

\item \textbf{Why?} Beta cannot produce the spread, which points to a
missing factor for value and mispricing for small growth. Within a size quintile
beta barely changes across BE/ME, and small growth has the highest beta of
all 25 portfolios despite the lowest return. Additional checks (not tabulated) show that the value spread moves as a common component of the CAPM residuals but does not lose money in market crashes, so it is not compensation for market risk. Small growth has
the highest skewness and idiosyncratic volatility of the 25 portfolios, a
negative alpha in every subsample, and stays mispriced even after adding a
value factor, which fits investors overpaying for lottery-like stocks. In both cases, the pattern lies outside what market beta can explain.

\end{itemize}

\subsection{Question (g)}

We replace RM$-$RF with the in-sample tangency portfolio of the 25
portfolios, with weights
\[
w_T = \frac{\hat\Sigma^{-1}\hat\mu}{\mathbf{1}'\hat\Sigma^{-1}\hat\mu}
\]
estimated over the full sample, and rerun the 25 time-series regressions and
the GRS test.

\begin{table}[H]
\centering
\caption{GRS test with the in-sample tangency portfolio of the 25 size and
BE/ME portfolios as the market proxy, July 1926 -- June 2026.}
\small
\begin{tabular}{lrr}
\toprule
& RM$-$RF & In-sample tangency \\
\midrule
GRS $F$-statistic & 3.37 & 0.00 \\
$p$-value & $<0.001$ & 1.00 \\
Largest $|\hat\alpha_i|$ (\% per month) & 0.53 & 0.00 \\
Proxy Sharpe ratio (monthly) & 0.131 & 0.299 \\
\bottomrule
\end{tabular}
\label{tab:p3g}
\vspace{2pt}\begin{minipage}{0.93\textwidth}\footnotesize
Tangency weights sum to one, with gross exposure $13.86$ and individual
weights between $-1.44$ and $1.06$. All 25 intercepts are below $10^{-14}$.
\end{minipage}
\end{table}

\paragraph{Discussion.}
\begin{itemize}
\setlength{\itemsep}{4pt}\setlength{\parskip}{0pt}

\item \textbf{Does the GRS test reject or fail to reject when we use the
in-sample tangency portfolio as the market proxy?} The GRS test fails to
reject. The statistic is zero and the
$p$-value is one: every intercept is zero up to rounding error, so the
in-sample tangency portfolio prices all 25 portfolios perfectly.

\item \textbf{Is this result surprising?} The result is not surprising.
The in-sample tangency portfolio is built from the same sample moments used
in the test, so it gives every portfolio a zero intercept by construction.
Since $w_T \propto \hat\Sigma^{-1}\hat\mu$, each portfolio's covariance with
the tangency portfolio is proportional to its mean, so
$\hat\mu_i = \hat\beta_{iT}\,\hat\mu_T$ and
$\hat\alpha_i = \hat\mu_i - \hat\beta_{iT}\hat\mu_T = 0$ for every $i$. In Sharpe-ratio terms, the proxy is already the
maximum-Sharpe-ratio combination of the test assets, so adding them cannot
improve it. The result therefore says nothing about the CAPM. The proxy is a
levered long--short portfolio of the test assets, not a market portfolio,
and it illustrates Roll's critique just as the perfect cross-sectional fit
did in Problem Set 3, Question (2).

\end{itemize}

\subsection{Question (h)}

Sample 1 contains odd-numbered months in even years and even-numbered months
in odd years; Sample 2 contains the rest, giving 600 months each. We estimate
tangency weights from Sample 1 moments and apply them to Sample 2 returns,
then reverse the roles. The combined, chronologically sorted series is the
out-of-sample tangency portfolio, which we use as the market proxy in the 25
time-series regressions.

\begin{table}[H]
\centering
\caption{GRS tests on the 25 size and BE/ME portfolios under three market
proxies, July 1926 -- June 2026.}
\small
\begin{tabular}{lrrr}
\toprule
& RM$-$RF & In-sample tangency & Out-of-sample tangency \\
\midrule
GRS $F$-statistic & 3.37 & 0.00 & 2.78 \\
$p$-value & $<0.001$ & 1.00 & $<0.001$ \\
Proxy Sharpe ratio (monthly) & 0.131 & 0.299 & 0.216 \\
\bottomrule
\end{tabular}
\label{tab:p3h}
\vspace{2pt}\begin{minipage}{0.93\textwidth}\footnotesize
Out-of-sample weights sum to one in each half, with gross exposure $14.81$
(estimated on Sample 1) and $18.44$ (Sample 2). The correlation between the
two weight vectors is $0.49$. All 25 intercepts are positive, between $0.14$
and $0.46\%$ per month.
\end{minipage}
\end{table}

\paragraph{Discussion.}
\begin{itemize}
\setlength{\itemsep}{4pt}\setlength{\parskip}{0pt}

\item \textbf{Does the GRS test reject or fail to reject when we use the
out-of-sample tangency portfolio as the market proxy?} The GRS test rejects.
The out-of-sample tangency portfolio
leaves significant pricing errors in the 25 portfolios.

\item \textbf{Why?} It rejects because the weights are estimated with
error.
The identity in (g) holds only when the weights and the test use the same
sample moments. Here they do not, and the mean estimates are noisy enough
that the two halves produce quite different weights. A portfolio that is
efficient in its own half loses about a third of its Sharpe ratio in the
other, so it lies inside the frontier of the test assets. Because the proxy is noisy, the 25 portfolios have low betas on it, so their beta-implied returns fall short of their average returns and all 25 intercepts are positive. The rejection therefore reflects the noisy proxy as much as the pricing model.

\item \textbf{Is this result surprising?} The result is not surprising.
Optimized weights fit the sample they come from and travel poorly, so once the weights come from a different sample, nothing forces the intercepts to zero and the test regains its power. The portfolio still beats the market's Sharpe ratio and
lowers the GRS statistic, but the test asks whether the proxy is efficient,
not whether it beats the market, and out of sample it is not.

\end{itemize}

\paragraph{AI integration: Roll's critique.} Before comparing (g) and (h),
we asked Claude Sonnet 5.5 (chat mode) why the in-sample tangency portfolio
makes the GRS test fail to reject trivially, and why the out-of-sample
version does not. It answered that GRS is a monotone function of
$(SR^{*2} - SR_f^2)/(1 + SR_f^2)$, that an in-sample tangency portfolio
attains $SR_f = SR^*$ by construction, and that weights estimated on a
different sample leave the proxy inside the test-sample frontier, restoring
power. It added that estimation error in those weights makes the test
over-reject as a test of the underlying model.

\begin{itemize}
\setlength{\itemsep}{4pt}\setlength{\parskip}{0pt}

\item[(i)] \textbf{Did the in-sample tangency portfolio in (g) produce a
GRS statistic near zero and a very high $p$-value? Did the LLM's explanation
correctly predict this?} Yes, on both counts. The in-sample tangency portfolio
produced a GRS statistic of zero and a $p$-value of one, exactly as the LLM
predicted.

\item[(ii)] \textbf{Does the LLM correctly identify that the in-sample
tangency portfolio makes all alphas exactly zero, or does it give a vaguer
explanation about overfitting?} The LLM identified the exact identity. It
showed that
$\mathrm{Cov}(R_i, f) = \hat\mu_i$, so $\hat\beta_i = \hat\mu_i/\mu_f$ and
every $\hat\alpha_i$ is zero, rather than appealing loosely to overfitting.

\item[(iii)] \textbf{Does the out-of-sample tangency portfolio in (h)
restore the GRS test's power? How does this reinforce the Roll critique?}
Yes, and the contrast reinforces Roll's critique.
With the same assets, the same data and the same model, the verdict moves
from a perfect fit to a clear rejection depending only on how the proxy is
built. The LLM's own caveat makes the same point without naming Roll: the
out-of-sample rejection partly reflects estimation error in the proxy, so
neither result is a clean verdict on the CAPM. The LLM did overstate one point: it claimed the exact $F$ distribution applies out of sample because the proxy is a given regressor. In our construction each half's proxy uses weights estimated from the other half's test-asset returns, so the $F$ distribution is only approximate. GRS can only ever test
whether the chosen proxy is efficient.

\end{itemize}

\subsection{Question (i)}

We replace RM$-$RF with the in-sample tangency portfolio of the 30
value-weight industry portfolios, estimated over the full sample, and rerun
the GRS test on the 25 size and BE/ME portfolios.

\begin{table}[H]
\centering
\caption{CAPM intercepts with the in-sample industry tangency portfolio as
the market proxy, percent per month.}
\small
\begin{tabular}{lrrrrrcrrrrr}
\toprule
& \multicolumn{5}{c}{$\hat\alpha_i$} && \multicolumn{5}{c}{$t(\hat\alpha_i)$} \\
\cmidrule{2-6}\cmidrule{8-12}
& L & 2 & 3 & 4 & H && L & 2 & 3 & 4 & H \\
\midrule
S & $-0.46$ & $-0.10$ & 0.26 & 0.47 & 0.62 && $-1.40$ & $-0.36$ & 1.06 & 2.03 & 2.44 \\
2 & $-0.17$ & 0.19 & 0.25 & 0.33 & 0.46 && $-0.81$ & 0.95 & 1.32 & 1.63 & 1.96 \\
3 & $-0.09$ & 0.14 & 0.19 & 0.27 & 0.31 && $-0.48$ & 0.86 & 1.14 & 1.46 & 1.37 \\
4 & $-0.10$ & 0.02 & 0.11 & 0.22 & 0.21 && $-0.64$ & 0.11 & 0.66 & 1.21 & 0.91 \\
B & $-0.12$ & $-0.07$ & 0.03 & $-0.05$ & 0.20 && $-1.00$ & $-0.54$ & 0.22 & $-0.31$ & 0.86 \\
\bottomrule
\end{tabular}
\label{tab:p3i}
\vspace{2pt}\begin{minipage}{0.93\textwidth}\footnotesize
GRS $= 3.51$, $p < 0.001$, against $3.37$ with RM$-$RF. Monthly Sharpe
ratio of the industry tangency portfolio: $0.236$.
\end{minipage}
\end{table}

\paragraph{Discussion.}
\begin{itemize}
\setlength{\itemsep}{4pt}\setlength{\parskip}{0pt}

\item \textbf{How well do we price the 25 size and BE/ME portfolios?}
The industry tangency portfolio prices the 25 portfolios no better than the
market. The GRS test still rejects, with a slightly larger
statistic than under RM$-$RF, and the intercepts keep the same pattern:
negative for small growth, positive and increasing in BE/ME elsewhere.

\item \textbf{Is this surprising? Why or why not?} This is not
surprising. The tangency portfolio is efficient
only relative to the assets it is built from. Industries are not sorted on
size or BE/ME, so their optimal combination carries little of the value
spread that drives the 25 portfolios' returns. A higher Sharpe ratio than
the market does not help: what matters is whether the proxy lies on the
frontier of the test assets, and the low correlation between the two tangency portfolios in (k) shows that it does not. Efficiency within one set of assets does not carry over to another.

\end{itemize}

\subsection{Question (j)}

We replace RM$-$RF with the in-sample tangency portfolio of the 10 past
return portfolios, estimated over the full sample, and rerun the GRS test on
the 25 size and BE/ME portfolios. Because the past return portfolios begin in
January 1927, this test uses 1927-01 to 2026-06 ($T = 1194$).

\begin{table}[H]
\centering
\caption{CAPM intercepts with the in-sample past-return tangency portfolio
as the market proxy, percent per month, 1927-01 to 2026-06.}
\small
\begin{tabular}{lrrrrrcrrrrr}
\toprule
& \multicolumn{5}{c}{$\hat\alpha_i$} && \multicolumn{5}{c}{$t(\hat\alpha_i)$} \\
\cmidrule{2-6}\cmidrule{8-12}
& L & 2 & 3 & 4 & H && L & 2 & 3 & 4 & H \\
\midrule
S & $-0.20$ & $-0.07$ & 0.35 & 0.50 & 0.74 && $-0.58$ & $-0.24$ & 1.39 & 2.10 & 2.77 \\
2 & $-0.17$ & 0.21 & 0.26 & 0.38 & 0.53 && $-0.77$ & 0.99 & 1.30 & 1.83 & 2.15 \\
3 & $-0.07$ & 0.12 & 0.19 & 0.31 & 0.39 && $-0.37$ & 0.71 & 1.09 & 1.60 & 1.65 \\
4 & $-0.11$ & 0.02 & 0.12 & 0.28 & 0.28 && $-0.66$ & 0.15 & 0.71 & 1.49 & 1.17 \\
B & $-0.07$ & $-0.07$ & 0.07 & $-0.00$ & 0.28 && $-0.54$ & $-0.52$ & 0.44 & $-0.02$ & 1.17 \\
\bottomrule
\end{tabular}
\label{tab:p3j}
\vspace{2pt}\begin{minipage}{0.93\textwidth}\footnotesize
GRS $= 3.36$, $p < 0.001$, $F(25, 1168)$, against $3.37$ with RM$-$RF over
the same sample. The tangency portfolio is mainly long Winner ($0.89$) and
short Loser ($-0.96$), with a monthly Sharpe ratio of $0.264$.
\end{minipage}
\end{table}

\paragraph{Discussion.}
\begin{itemize}
\setlength{\itemsep}{4pt}\setlength{\parskip}{0pt}

\item \textbf{How well do we price the 25 size and BE/ME portfolios?}
No better than the market. The GRS statistic is essentially unchanged from
the CAPM, but the intercepts now rise even more steeply with BE/ME, reaching
$0.74\%$ per month for small value, and their average absolute size
increases by about half.

\item \textbf{Is this surprising? Why or why not?} The result is surprising
at first, but understandable. We expected the momentum tangency portfolio to
do better, since its Sharpe ratio is about twice the market's. As in (i),
however, the test asks whether the proxy is efficient relative
to the 25 portfolios, and a bet on winners over losers carries no
information about value. It even works against value: value portfolios move
with past losers, so the proxy expects little from them and their alphas
grow. The GRS statistic holds steady only because those larger alphas are
also noisier. A better portfolio is not necessarily a better proxy.

\end{itemize}

\subsection{Question (k)}

We form the in-sample tangency portfolio of each of the three portfolio sets
and correlate their monthly returns over the common sample, 1927-01 to
2026-06 ($T = 1194$).

\begin{table}[H]
\centering
\caption{Correlation matrix of the in-sample tangency portfolios, 1927-01 to
2026-06.}
\small
\begin{tabular}{lrrr}
\toprule
& 30 industries & 10 past return & 25 size and BE/ME \\
\midrule
30 industries & 1.00 & 0.40 & 0.27 \\
10 past return & 0.40 & 1.00 & 0.25 \\
25 size and BE/ME & 0.27 & 0.25 & 1.00 \\
\bottomrule
\end{tabular}
\label{tab:p3k}
\vspace{2pt}\begin{minipage}{0.93\textwidth}\footnotesize
Monthly Sharpe ratios: $0.232$, $0.264$ and $0.299$. Correlations with
RM$-$RF: $0.61$, $0.52$ and $0.44$.
\end{minipage}
\end{table}

\paragraph{Discussion.}
\begin{itemize}
\setlength{\itemsep}{4pt}\setlength{\parskip}{0pt}

\item \textbf{Comment on these correlations in light of (i) and (j).} The
three tangency portfolios are only weakly correlated, so they are driven by
separate factors rather than one beta measured three ways. Industry
differences, past winners over losers and value are different sources of
return. This explains (i) and (j): a proxy that shares little of the 25 portfolios' movement cannot lie on their frontier. A single factor
that prices all three sets would have to span all three directions at once.

\end{itemize}


% =====================================================================
\section{Capstone Exercise: Choose a Factor, Build a Test}
% =====================================================================

\subsection{Step 1}

Choose RMW factor.

\subsection{Step 2}

We use the 25 size and BE/ME portfolios as the test assets and RMW (Robust
Minus Weak) as the additional factor. We expect RMW to help price these
portfolios because profitability captures a dimension of expected returns
that is not included in the CAPM market factor and may vary systematically
across firms with different size and book-to-market characteristics.

Before examining the results, our hypothesis is that adding RMW to the market
factor will reduce the GRS F-statistic relative to the CAPM, although we do
not expect it to eliminate all pricing errors. We expect the largest alpha
reductions to occur for portfolios with stronger RMW exposures, while
portfolios with weak profitability exposure should show smaller
improvements.

Economically, RMW may earn a positive risk premium because firms with weak
profitability could be riskier and therefore require different expected
returns, or the profitability premium may reflect systematic mispricing if
investors do not fully incorporate information about firms' profitability
into prices. The test will examine whether adding RMW materially reduces the
pricing errors of the 25 portfolios beyond what market beta alone can
explain.

\subsection{Question 3a}

The notebook cells are reproduced in order. The GRS validation output
follows the validation cell, and the full alpha comparison follows the last
cell.

\begin{lstlisting}
import pandas as pd
import numpy as np
import statsmodels.api as sm
from scipy.stats import f

excel_file = "Problem_Set4_2026.xlsx"
factor_file = "F-F_Research_Data_5_Factors_2x3.csv"
pd.set_option("display.max_columns", 35)
pd.set_option("display.width", 200)

book = pd.ExcelFile(excel_file)
print("Excel sheets:", book.sheet_names)

portfolio_raw = pd.read_excel(
    book, sheet_name="25 size and BEME portfolios", header=None
)
market_raw = pd.read_excel(
    book, sheet_name="Market, Rf", header=None
)

print("\nPortfolio header rows:")
print(portfolio_raw.iloc[:3].to_string(index=False, header=False))

print("\nMarket header rows:")
print(market_raw.iloc[:2].to_string(index=False, header=False))

# Preserve the workbook's original size-first, BE/ME-second order.
portfolio_cols = [
    f"S{s}_BM{b}" for s in range(1, 6) for b in range(1, 6)
]

portfolios = portfolio_raw.iloc[3:].copy()
assert portfolios.shape[1] == 26
portfolios.columns = ["Date"] + portfolio_cols

market = market_raw.iloc[2:].copy()
assert market.shape[1] == 3
assert market_raw.iloc[1, 1:].tolist() == ["RM-RF", "RF"]
market.columns = ["Date", "RM-RF", "RF"]

# The CSV header is on line 5.
ff = pd.read_csv(factor_file, skiprows=4)
ff.columns = ff.columns.str.strip()
ff = ff.rename(columns={ff.columns[0]: "Date"})

print("\nCSV columns:", ff.columns.tolist())

# Monthly dates have six digits. Annual dates have four.
date_text = ff["Date"].astype(str).str.strip()
monthly_rows = date_text.str.fullmatch(r"\d{6}")

print("CSV rows outside monthly section:", int((~monthly_rows).sum()))

# Retain only RMW from the external factor file.
rmw = ff.loc[monthly_rows, ["Date", "RMW"]].copy()

for name, data in [
    ("Portfolios", portfolios),
    ("Workbook market and RF", market),
    ("Monthly RMW", rmw)
]:
    print("\n" + name)
    print("Dimensions:", data.shape)
    print("Columns:", data.columns.tolist())
    print(data.head().to_string(index=False))
    print("Data types:")
    print(data.dtypes)
\end{lstlisting}
\begin{lstlisting}
sources = [
    ("Portfolios", portfolios),
    ("Workbook market and RF", market),
    ("Monthly RMW", rmw)
]

for name, data in sources:
    dates = pd.to_numeric(data["Date"], errors="raise")

    assert dates.notna().all(), name + ": missing dates"
    assert (dates % 1 == 0).all(), name + ": invalid YYYYMM dates"

    data["Date"] = pd.to_datetime(
        dates.astype(int).astype(str), format="%Y%m", errors="raise"
    ).dt.to_period("M")

    value_cols = data.columns.drop("Date")
    data[value_cols] = data[value_cols].apply(
        pd.to_numeric, errors="raise"
    )

    # Recognize standard missing-return codes explicitly.
    sentinel_count = data[value_cols].isin([-99.99, -999]).sum().sum()
    print(f"\n{name}: missing-value codes = {sentinel_count}")

    data[value_cols] = data[value_cols].replace(
        [-99.99, -999], np.nan
    )

    print("Missing values before merging:")
    print(data[value_cols].isna().sum().to_string())
    print("Rows with any missing return:",
          int(data[value_cols].isna().any(axis=1).sum()))
    print("Duplicate dates:", int(data["Date"].duplicated().sum()))
    print("Range:", data["Date"].min(), "to", data["Date"].max())

    assert not data["Date"].duplicated().any(), name + ": duplicate dates"

# Both models will use this single merged sample.
common = portfolios.merge(
    market, on="Date", how="inner", validate="one_to_one"
)
common = common.merge(
    rmw, on="Date", how="inner", validate="one_to_one"
)
common = common.sort_values("Date").set_index("Date")

for name, data in sources:
    excluded = data.loc[~data["Date"].isin(common.index), "Date"]
    print(f"\n{name}: {len(excluded)} months excluded by date overlap")
    if len(excluded) > 0:
        print("Excluded range:", excluded.min(), "to", excluded.max())

print("\nMissing values after merging:")
print(common.isna().sum().to_string())
print("Rows with any missing value:",
      int(common.isna().any(axis=1).sum()))

assert len(common) > 0, "No overlapping observations."
assert not common.isna().any().any(), "Missing values remain. Stop here."
assert np.isfinite(common.to_numpy()).all(), "Nonfinite values remain."

expected_months = pd.period_range(
    common.index.min(), common.index.max(), freq="M"
)
assert common.index.equals(expected_months), "There are gaps in the sample."

print("\nCommon sample:")
print("Start:", common.index.min())
print("End:", common.index.max())
print("T:", len(common))
print("Units: percentage points for portfolios, RM-RF, RF, and RMW.")
print("\nCleaned data types:")
print(common.dtypes)
\end{lstlisting}
\begin{lstlisting}
Y = common[portfolio_cols].sub(common["RF"], axis=0)

F_capm = common[["RM-RF"]].copy()
F_aug = common[["RM-RF", "RMW"]].copy()

assert Y.index.equals(F_capm.index)
assert Y.index.equals(F_aug.index)
assert Y.shape[1] == 25

print("Portfolio excess returns = raw portfolio returns - workbook RF.")
print("RM-RF and RMW are used directly.")
print("All returns remain in percentage points.")
print(Y.head().to_string())
\end{lstlisting}
\begin{lstlisting}
X_capm = sm.add_constant(F_capm, has_constant="add")

capm_rows = []
capm_residuals = pd.DataFrame(index=Y.index)
capm_fitted = pd.DataFrame(index=Y.index)

for portfolio in portfolio_cols:
    model = sm.OLS(Y[portfolio], X_capm, missing="raise").fit()

    capm_rows.append({
        "Portfolio": portfolio,
        "Alpha": model.params["const"],
        "Alpha SE": model.bse["const"],
        "Alpha t-stat": model.tvalues["const"],
        "Market beta": model.params["RM-RF"],
        "Market beta t-stat": model.tvalues["RM-RF"]
    })

    capm_residuals[portfolio] = model.resid
    capm_fitted[portfolio] = model.fittedvalues

capm_results = pd.DataFrame(capm_rows).set_index("Portfolio")
print(capm_results.to_string(float_format=lambda x: f"{x:.6f}"))
\end{lstlisting}
\begin{lstlisting}
def grs_test(alpha, residuals, factors):
    alpha = np.asarray(alpha, dtype=float).reshape(-1)
    E = np.asarray(residuals, dtype=float)
    F = np.asarray(factors, dtype=float)

    assert E.ndim == 2 and F.ndim == 2

    T, N = E.shape
    K = F.shape[1]

    assert F.shape[0] == T
    assert len(alpha) == N
    assert T > N + K
    assert np.isfinite(alpha).all()
    assert np.isfinite(E).all()
    assert np.isfinite(F).all()

    # Covariance of OLS residuals, not raw portfolio returns.
    sigma = E.T @ E / (T - K - 1)

    factor_mean = F.mean(axis=0)
    centered_factors = F - factor_mean
    omega = centered_factors.T @ centered_factors / (T - 1)

    # Stop if covariance matrices are not positive definite.
    np.linalg.cholesky(sigma)
    np.linalg.cholesky(omega)

    alpha_term = alpha @ np.linalg.solve(sigma, alpha)
    factor_term = factor_mean @ np.linalg.solve(omega, factor_mean)

    statistic = ((T - N - K) / N) * alpha_term / (1 + factor_term)
    p_value = f.sf(statistic, N, T - N - K)

    return {
        "T": T,
        "N": N,
        "K": K,
        "GRS F-statistic": float(statistic),
        "p-value": float(p_value),
        "Residual denominator": T - K - 1,
        "Factor denominator": T - 1,
        "F df numerator": N,
        "F df denominator": T - N - K
    }

capm_grs = grs_test(
    capm_results["Alpha"], capm_residuals, F_capm
)

print(pd.Series(capm_grs).to_string())
\end{lstlisting}
\begin{lstlisting}
X_aug = sm.add_constant(F_aug, has_constant="add")

aug_rows = []
aug_residuals = pd.DataFrame(index=Y.index)
aug_fitted = pd.DataFrame(index=Y.index)

for portfolio in portfolio_cols:
    model = sm.OLS(Y[portfolio], X_aug, missing="raise").fit()

    aug_rows.append({
        "Portfolio": portfolio,
        "Alpha": model.params["const"],
        "Alpha SE": model.bse["const"],
        "Alpha t-stat": model.tvalues["const"],
        "Market beta": model.params["RM-RF"],
        "Market beta t-stat": model.tvalues["RM-RF"],
        "RMW beta": model.params["RMW"],
        "RMW beta t-stat": model.tvalues["RMW"]
    })

    aug_residuals[portfolio] = model.resid
    aug_fitted[portfolio] = model.fittedvalues

aug_results = pd.DataFrame(aug_rows).set_index("Portfolio")
print(aug_results.to_string(float_format=lambda x: f"{x:.6f}"))
\end{lstlisting}
\begin{lstlisting}
aug_grs = grs_test(
    aug_results["Alpha"], aug_residuals, F_aug
)

print(pd.Series(aug_grs).to_string())
\end{lstlisting}
\begin{lstlisting}
validation_passed = False
validation_rows = []
audit_rows = []

models = [
    ("CAPM", capm_results, capm_residuals, capm_fitted,
     F_capm, capm_grs, 1),
    ("Market + RMW", aug_results, aug_residuals, aug_fitted,
     F_aug, aug_grs, 2)
]

for name, results, residuals, fitted, factors, result, expected_k in models:
    alpha = results["Alpha"].to_numpy()
    original_grs = result["GRS F-statistic"]

    # Check 1: zero alpha.
    zero_grs = grs_test(
        np.zeros(len(alpha)), residuals, factors
    )["GRS F-statistic"]

    # Check 2: reverse the portfolio order consistently.
    reorder_grs = grs_test(
        alpha[::-1], residuals.iloc[:, ::-1], factors
    )["GRS F-statistic"]

    # Check 3: scale all returns and factors, then refit.
    scaled_X = sm.add_constant(factors * 100, has_constant="add")
    scaled_alpha = []
    scaled_residuals = pd.DataFrame(index=Y.index)

    for portfolio in portfolio_cols:
        scaled_model = sm.OLS(Y[portfolio] * 100, scaled_X).fit()
        scaled_alpha.append(scaled_model.params["const"])
        scaled_residuals[portfolio] = scaled_model.resid

    scale_grs = grs_test(
        scaled_alpha, scaled_residuals, factors * 100
    )["GRS F-statistic"]

    checks = [
        ("Zero alpha", zero_grs,
         np.isclose(zero_grs, 0, rtol=0, atol=1e-12)),
        ("Reorder invariance", reorder_grs,
         np.isclose(reorder_grs, original_grs, rtol=1e-9, atol=1e-12)),
        ("Scale invariance", scale_grs,
         np.isclose(scale_grs, original_grs, rtol=1e-9, atol=1e-12))
    ]

    for check, value, passed in checks:
        validation_rows.append({
            "Model": name,
            "Check": check,
            "GRS": value,
            "Result": "PASS" if passed else "FAIL"
        })

    T, N = residuals.shape
    K = factors.shape[1]

    # Verify residuals equal observed excess returns minus fitted returns.
    residual_check = np.allclose(
        residuals.to_numpy(),
        (Y - fitted).to_numpy(),
        rtol=1e-10, atol=1e-10
    )

    # OLS residuals should be orthogonal to the regressors.
    X = sm.add_constant(factors, has_constant="add").to_numpy()
    orthogonal_check = np.allclose(
        X.T @ residuals.to_numpy(), 0, rtol=0, atol=1e-7
    )

    specification_check = (
        K == expected_k
        and result["K"] == expected_k
        and result["N"] == 25
        and result["Residual denominator"] == T - K - 1
        and result["Factor denominator"] == T - 1
        and result["F df numerator"] == N
        and result["F df denominator"] == T - N - K
    )

    audit_rows.append({
        "Model": name,
        "K": K,
        "Residual denominator": result["Residual denominator"],
        "Factor denominator": result["Factor denominator"],
        "F degrees of freedom": (N, T - N - K),
        "OLS residual check": "PASS" if residual_check else "FAIL",
        "Orthogonality": "PASS" if orthogonal_check else "FAIL",
        "Specification": "PASS" if specification_check else "FAIL"
    })

validation_results = pd.DataFrame(validation_rows)
grs_audit = pd.DataFrame(audit_rows)

print(validation_results.to_string(index=False))
print("\nGRS implementation audit:")
print(grs_audit.to_string(index=False))

assert validation_results["Result"].eq("PASS").all(), "GRS validation failed."
assert grs_audit["OLS residual check"].eq("PASS").all(), "Residual check failed."
assert grs_audit["Orthogonality"].eq("PASS").all(), "OLS orthogonality failed."
assert grs_audit["Specification"].eq("PASS").all(), "GRS specification failed."

validation_passed = True
print("\nAll checks passed. Continue to the comparison tables.")
\end{lstlisting}

\begin{table}[H]
\centering
\caption{GRS validation, capstone implementation, 1963-07 to 2026-06
($T = 756$).}
\small
\begin{tabular}{llrl}
\toprule
Model & Check & GRS & Result \\
\midrule
CAPM & Zero alpha & 0.000000 & PASS \\
CAPM & Reorder invariance & 4.164336 & PASS \\
CAPM & Scale invariance & 4.164336 & PASS \\
Market + RMW & Zero alpha & 0.000000 & PASS \\
Market + RMW & Reorder invariance & 3.723390 & PASS \\
Market + RMW & Scale invariance & 3.723390 & PASS \\
\bottomrule
\end{tabular}
\label{tab:cap_valid}
\vspace{2pt}\begin{minipage}{0.93\textwidth}\footnotesize
Residual covariance denominators: 754 (CAPM) and 753 (Market + RMW); factor
covariance denominator 755; $F(25, 730)$ and $F(25, 729)$.
\end{minipage}
\end{table}

\begin{lstlisting}
assert validation_passed, "Run the validation cell first."

comparison = pd.DataFrame([
    {
        "Model": "CAPM",
        "Factors": "RM-RF",
        "T": capm_grs["T"],
        "N": capm_grs["N"],
        "K": capm_grs["K"],
        "GRS F-statistic": capm_grs["GRS F-statistic"],
        "p-value": capm_grs["p-value"],
        "Average RMW return": np.nan
    },
    {
        "Model": "Augmented model",
        "Factors": "RM-RF + RMW",
        "T": aug_grs["T"],
        "N": aug_grs["N"],
        "K": aug_grs["K"],
        "GRS F-statistic": aug_grs["GRS F-statistic"],
        "p-value": aug_grs["p-value"],
        "Average RMW return": common["RMW"].mean()
    }
])

assert capm_grs["GRS F-statistic"] > 0

grs_reduction = (
    (capm_grs["GRS F-statistic"] - aug_grs["GRS F-statistic"])
    / capm_grs["GRS F-statistic"] * 100
)

print(comparison.to_string(
    index=False, na_rep="N/A",
    formatters={
        "GRS F-statistic": "{:.6f}".format,
        "p-value": "{:.6e}".format,
        "Average RMW return": "{:.6f}".format
    }
))
print(f"\nGRS percentage reduction: {grs_reduction:.6f}%")
\end{lstlisting}
\begin{lstlisting}
assert validation_passed, "Run the validation cell first."

alpha_comparison = pd.DataFrame({
    "Portfolio": portfolio_cols,
    "CAPM alpha": capm_results["Alpha"].to_numpy(),
    "CAPM alpha t-stat": capm_results["Alpha t-stat"].to_numpy(),
    "Market + RMW alpha": aug_results["Alpha"].to_numpy(),
    "Market + RMW alpha t-stat": aug_results["Alpha t-stat"].to_numpy(),
    "RMW beta": aug_results["RMW beta"].to_numpy(),
    "RMW beta t-stat": aug_results["RMW beta t-stat"].to_numpy()
})

alpha_comparison["Absolute alpha reduction"] = (
    alpha_comparison["CAPM alpha"].abs()
    - alpha_comparison["Market + RMW alpha"].abs()
)

alpha_sorted = alpha_comparison.sort_values(
    "Absolute alpha reduction", ascending=False
).copy()

print("Original portfolio order:")
print(alpha_comparison.to_string(
    index=False, float_format=lambda x: f"{x:.6f}"
))

print("\nSorted by absolute alpha reduction:")
print(alpha_sorted.to_string(
    index=False, float_format=lambda x: f"{x:.6f}"
))
\end{lstlisting}
\begin{lstlisting}
assert validation_passed, "Run the validation cell first."

T = len(common)
rmw_mean = common["RMW"].mean()
rmw_sd = common["RMW"].std(ddof=1)
rmw_se = rmw_sd / np.sqrt(T)

assert rmw_se > 0
rmw_t = rmw_mean / rmw_se

rmw_mean_test = pd.DataFrame([{
    "T": T,
    "Mean monthly RMW (%)": rmw_mean,
    "Standard deviation (%)": rmw_sd,
    "Standard error (%)": rmw_se,
    "Mean t-statistic": rmw_t
}])

print("One-sample test: H0: E[RMW] = 0")
print(rmw_mean_test.to_string(
    index=False, float_format=lambda x: f"{x:.6f}"
))
print("\nThis tests the RMW mean, not the joint portfolio pricing errors.")
\end{lstlisting}
\begin{lstlisting}
from scipy.stats import t as t_dist

return_sets = {
    "CAPM factors": F_capm,
    "Augmented factors": F_aug,
    "CAPM factors + portfolios": pd.concat([F_capm, Y], axis=1),
    "Augmented factors + portfolios": pd.concat([F_aug, Y], axis=1)
}

sr2 = {}

for name, returns in return_sets.items():
    mean = returns.mean().to_numpy()
    covariance = returns.cov().to_numpy()
    sr2[name] = float(mean @ np.linalg.solve(covariance, mean))

capm_gain = sr2["CAPM factors + portfolios"] - sr2["CAPM factors"]
aug_gain = sr2["Augmented factors + portfolios"] - sr2["Augmented factors"]
gain_reduction = (1 - aug_gain / capm_gain) * 100

print("SHARPE-RATIO RESULTS (monthly)")
for name, value in sr2.items():
    print(f"{name}: SR = {np.sqrt(value):.4f}, SR² = {value:.6f}")

print(f"\nCAPM incremental SR²: {capm_gain:.6f}")
print(f"Augmented incremental SR²: {aug_gain:.6f}")
print(f"Incremental SR² reduction: {gain_reduction:.6f}%")
print(f"CAPM scaling denominator: {1 + sr2['CAPM factors']:.4f}")
print(f"Augmented scaling denominator: {1 + sr2['Augmented factors']:.4f}")

rmw_p = 2 * t_dist.sf(abs(rmw_t), df=len(common) - 1)

print("\nFACTOR MEANS AND RMW MEAN TEST")
print(f"Mean RM-RF: {common['RM-RF'].mean():.6f}% per month")
print(f"Mean RMW: {rmw_mean:.6f}% per month")
print(f"RMW mean t-statistic: {rmw_t:.6f}")
print(f"RMW mean-test p-value: {rmw_p:.6f}")

capm_abs = capm_results["Alpha"].abs()
aug_abs = aug_results["Alpha"].abs()
alpha_reduction = capm_abs - aug_abs

print("\nABSOLUTE-ALPHA SUMMARY (monthly percentage points)")
print(f"CAPM total absolute alpha: {capm_abs.sum():.6f}")
print(f"Augmented total absolute alpha: {aug_abs.sum():.6f}")
print(f"Portfolios with lower absolute alpha: {(alpha_reduction > 0).sum()} of 25")
\end{lstlisting}
\begin{lstlisting}
assert validation_passed, "Run the validation cell first."

print("COMMON SAMPLE")
print("Start:", common.index.min())
print("End:", common.index.max())
print("T:", len(common))

print("\nGRS RESULTS")
print(
    f"CAPM: F = {capm_grs['GRS F-statistic']:.6f}, "
    f"p = {capm_grs['p-value']:.6e}"
)
print(
    f"Market + RMW: F = {aug_grs['GRS F-statistic']:.6f}, "
    f"p = {aug_grs['p-value']:.6e}"
)
print(f"GRS percentage reduction: {grs_reduction:.6f}%")

print("\nRMW MEAN TEST")
print(rmw_mean_test.to_string(
    index=False, float_format=lambda x: f"{x:.6f}"
))

print("\nMODEL COMPARISON")
print(comparison.to_string(
    index=False, na_rep="N/A",
    formatters={
        "GRS F-statistic": "{:.6f}".format,
        "p-value": "{:.6e}".format,
        "Average RMW return": "{:.6f}".format
    }
))

print("\nALPHA COMPARISON: ORIGINAL ORDER")
print(alpha_comparison.to_string(
    index=False, float_format=lambda x: f"{x:.6f}"
))

print("\nALPHA COMPARISON: SORTED COPY")
print(alpha_sorted.to_string(
    index=False, float_format=lambda x: f"{x:.6f}"
))

print("\nGRS VALIDATION")
print(validation_results.to_string(index=False))
print(grs_audit.to_string(index=False))

print("\nTECHNICAL NOTE")
print(
    "The implementation uses one common sample, workbook RM-RF and RF, "
    "and only RMW from the external CSV. RF is subtracted once from "
    "raw portfolio returns. GRS follows the assignment's displayed "
    "formula, using OLS residual covariance divided by T-K-1, factor "
    "covariance divided by T-1, and F(N, T-N-K) degrees of freedom."
)
print(
    "Special attention was required for the workbook's multiple header "
    "rows, YYYYMM dates, and the CSV's separate annual section. "
    "The portfolio labels preserve the original workbook ordering."
)
print(
    "Date alignment excludes 444 workbook months before July 1963 "
    "and two RMW months after June 2026 in these uploaded files. "
    "No missing observations remain in the 756-month common sample. "
)
\end{lstlisting}
\begin{lstlisting}
pd.set_option("display.max_rows", None)
pd.set_option("display.max_columns", None)
pd.set_option("display.width", None)
pd.set_option("display.max_colwidth", None)

print(alpha_comparison.to_string(
    index=False,
    float_format=lambda x: f"{x:.6f}"
))
\end{lstlisting}

\begin{table}[H]
\centering
\caption{Alpha comparison, CAPM vs.\ Market + RMW, 25 size and BE/ME
portfolios, monthly percentage points. S1 = smallest size quintile, BM1 =
lowest BE/ME.}
\footnotesize
\setlength{\tabcolsep}{4pt}
\begin{tabular}{lrrrrrrr}
\toprule
& \multicolumn{2}{c}{CAPM} & \multicolumn{2}{c}{Market + RMW} & \multicolumn{2}{c}{RMW loading} & \\
\cmidrule(lr){2-3}\cmidrule(lr){4-5}\cmidrule(lr){6-7}
Portfolio & $\hat\alpha$ & $t$ & $\hat\alpha$ & $t$ & $\hat\beta_{RMW}$ & $t$ & $|\hat\alpha_{CAPM}| - |\hat\alpha_{aug}|$ \\
\midrule
\texttt{S1\_BM1} & $-$0.522 & $-$2.93 & $-$0.212 & $-$1.33 & $-$1.029 & $-$14.55 & 0.310 \\
\texttt{S1\_BM2} & 0.059 & 0.38 & 0.328 & 2.32 & $-$0.891 & $-$14.17 & $-$0.268 \\
\texttt{S1\_BM3} & 0.137 & 1.06 & 0.296 & 2.40 & $-$0.527 & $-$9.60 & $-$0.159 \\
\texttt{S1\_BM4} & 0.347 & 2.64 & 0.487 & 3.84 & $-$0.465 & $-$8.23 & $-$0.140 \\
\texttt{S1\_BM5} & 0.473 & 3.16 & 0.598 & 4.06 & $-$0.416 & $-$6.33 & $-$0.125 \\
\texttt{S2\_BM1} & $-$0.266 & $-$1.96 & $-$0.070 & $-$0.55 & $-$0.651 & $-$11.58 & 0.196 \\
\texttt{S2\_BM2} & 0.103 & 0.92 & 0.213 & 1.96 & $-$0.368 & $-$7.59 & $-$0.111 \\
\texttt{S2\_BM3} & 0.245 & 2.37 & 0.285 & 2.74 & $-$0.132 & $-$2.85 & $-$0.040 \\
\texttt{S2\_BM4} & 0.291 & 2.75 & 0.344 & 3.24 & $-$0.174 & $-$3.68 & $-$0.052 \\
\texttt{S2\_BM5} & 0.320 & 2.43 & 0.405 & 3.10 & $-$0.283 & $-$4.86 & $-$0.085 \\
\texttt{S3\_BM1} & $-$0.220 & $-$2.01 & $-$0.089 & $-$0.85 & $-$0.434 & $-$9.34 & 0.131 \\
\texttt{S3\_BM2} & 0.124 & 1.43 & 0.158 & 1.81 & $-$0.110 & $-$2.83 & $-$0.033 \\
\texttt{S3\_BM3} & 0.144 & 1.69 & 0.140 & 1.62 & 0.014 & 0.36 & 0.004 \\
\texttt{S3\_BM4} & 0.272 & 2.85 & 0.276 & 2.86 & $-$0.013 & $-$0.29 & $-$0.004 \\
\texttt{S3\_BM5} & 0.344 & 2.76 & 0.359 & 2.85 & $-$0.050 & $-$0.89 & $-$0.015 \\
\texttt{S4\_BM1} & $-$0.069 & $-$0.83 & 0.020 & 0.25 & $-$0.297 & $-$8.24 & 0.049 \\
\texttt{S4\_BM2} & 0.016 & 0.23 & $-$0.002 & $-$0.02 & 0.058 & 1.90 & 0.014 \\
\texttt{S4\_BM3} & 0.143 & 1.81 & 0.109 & 1.38 & 0.113 & 3.20 & 0.034 \\
\texttt{S4\_BM4} & 0.276 & 3.04 & 0.270 & 2.94 & 0.022 & 0.53 & 0.007 \\
\texttt{S4\_BM5} & 0.221 & 1.84 & 0.236 & 1.95 & $-$0.049 & $-$0.92 & $-$0.015 \\
\texttt{S5\_BM1} & 0.020 & 0.33 & $-$0.049 & $-$0.87 & 0.229 & 9.11 & $-$0.030 \\
\texttt{S5\_BM2} & 0.030 & 0.51 & $-$0.031 & $-$0.56 & 0.201 & 8.09 & $-$0.002 \\
\texttt{S5\_BM3} & 0.091 & 1.23 & 0.054 & 0.72 & 0.125 & 3.77 & 0.038 \\
\texttt{S5\_BM4} & 0.021 & 0.22 & $-$0.025 & $-$0.26 & 0.152 & 3.57 & $-$0.004 \\
\texttt{S5\_BM5} & 0.194 & 1.47 & 0.234 & 1.76 & $-$0.133 & $-$2.24 & $-$0.040 \\
\bottomrule
\end{tabular}
\label{tab:cap_alpha}
\end{table}

\subsection{3a report}

\paragraph{What the code got right.} The code correctly constructed
portfolio excess returns, used the same 1963-07 to 2026-06 common sample for
both models, and estimated the CAPM with ($K=1$) and the augmented model with
($K=2$). The GRS implementation used the residual covariance matrix with
denominator ($T-K-1$), the factor covariance matrix with denominator ($T-1$),
the correct scaling term, and ($F(N,T-N-K)$) degrees of freedom. The
zero-alpha, reorder-invariance, and scale-invariance checks also passed, and
the independently recomputed GRS statistics matched the notebook.

\paragraph{What we corrected or added.} We did not need to correct the
regression or GRS calculations themselves. The main correction was to the
alpha-table display: the original notebook output was visually truncated, so
we added a separate cell that prints all 25 portfolios for Step 3(b).
Additional Sharpe-ratio and summary statistics were added later for Step 4,
rather than to fix an error in the Step 3 implementation.

\subsection{Step 3b}

\paragraph{LLM's Explaination.} Adding RMW helps price a relatively
small subset of portfolios, especially the low-BE/ME portfolios. The largest
improvements are for S1\_BM1 (absolute alpha falls by 0.310), S2\_BM1 (0.196),
and S3\_BM1 (0.131). Smaller improvements occur for portfolios such as
S4\_BM1, S4\_BM3, and S5\_BM3. RMW does not help many of the remaining
portfolios. The largest deteriorations are S1\_BM2 ($-0.268$), S1\_BM3
($-0.159$), S1\_BM4 ($-0.140$), and S1\_BM5 ($-0.125$). Overall, RMW appears
most useful for some low-BE/ME portfolios, particularly among smaller firms,
but it does not consistently reduce pricing errors across the full set of 25
portfolios.

\paragraph{Citical Assesssing.} The LLM's portfolio-level figures are
accurate, and it correctly identifies the largest reductions (S1\_BM1,
S2\_BM1, S3\_BM1) and deteriorations (S1\_BM2--S1\_BM5) in absolute alpha.
However, it understates how limited the improvement is: absolute alphas fall
for only 9 of 25 portfolios, several of those reductions are below 0.015, and
total absolute alpha rises from 4.95 to 5.29. Its claim that RMW helps
``particularly among smaller firms'' is misleading, because the
smallest-size quintile also contains the four largest deteriorations. Most
importantly, it misses the pattern that organizes the table: given RMW's
positive sample mean, pricing errors shrink only where the sign of the RMW
loading offsets the sign of the CAPM alpha, and grow where the loading
pushes an already-positive alpha further from zero, as in small value
portfolios. The table therefore only partly supports our hypothesis: weakly
exposed portfolios change little, but strong RMW exposure is associated with
large alpha changes in either direction, not necessarily reductions.

\subsection{Step 4}

\begin{table}[H]
\centering
\caption{GRS tests, CAPM vs.\ augmented model, 25 size and BE/ME portfolios,
1963-07 to 2026-06 ($T = 756$).}
\small
\begin{tabular}{llrrrr}
\toprule
Model & Factors & GRS F-statistic & $p$-value & Mean RM-RF & Mean RMW \\
\midrule
CAPM & RM-RF & 4.164336 & $9.853287\times10^{-11}$ & 0.600794\% & --- \\
Augmented & RM-RF + RMW & 3.723390 & $4.213838\times10^{-9}$ & 0.600794\% & 0.242804\% \\
\bottomrule
\end{tabular}
\label{tab:cap_grs}
\end{table}

\paragraph{4(a)} Adding RMW lowers GRS from 4.164336 to 3.723390, a
10.588612\% reduction. This indicates improvement in the joint pricing-error
criterion, but the reduction itself is not a formal test of statistically
significant improvement. The augmented-model $p$-value remains
$4.213838\times10^{-9}$, strongly rejecting the joint null that all 25 alphas
equal zero.

Using the same uploaded sample and sample covariances with denominator
$T-1$, the incremental maximum squared monthly Sharpe ratio from adding the
test assets falls from 0.145390 to 0.132731, approximately 8.707282\%. Thus,
the test assets offer a smaller remaining Sharpe-ratio improvement, but no
combination of RM-RF and RMW is mean-variance efficient relative to the
factors plus the 25 test assets. Most of the change comes from the factor
side: the maximum monthly Sharpe ratio of the factors rises from 0.1346
(market alone) to 0.1917 (market + RMW), while that of the factors plus test
assets rises only from 0.4044 to 0.4117. The percentage differs from the GRS
reduction mainly because GRS divides this increment by $1+SR_f^2$, which
rises from 1.0181 to 1.0367 when RMW is added; the remainder reflects the
degrees-of-freedom and covariance-denominator adjustments.

\paragraph{4(b)} Defining improvement as
$|\alpha_{CAPM}|-|\alpha_{augmented}|$, the largest reductions occur in
S1\_BM1 (0.310156), S2\_BM1 (0.196118), and S3\_BM1 (0.130670), measured in
monthly percentage points. These are the lowest-BE/ME portfolios in the first
three size quintiles. However, absolute alphas decline for only 9 of 25
portfolios. The largest deteriorations occur in S1\_BM2 ($-0.268286$),
S1\_BM3 ($-0.158904$), S1\_BM4 ($-0.140086$), and S1\_BM5 ($-0.125237$).
Stronger absolute RMW exposure therefore does not consistently imply
improvement: S1\_BM2 has a large loading of $-0.890504$, yet its positive
alpha increases substantially.

The hypothesis is partially supported. GRS falls and pricing errors remain,
as predicted, and the three largest reductions occur in portfolios with
strong negative RMW loadings ($-1.029$, $-0.651$, $-0.434$). Weakly exposed
portfolios ($|\beta_{RMW}| < 0.06$) change by at most 0.015, also consistent
with the hypothesis. However, strong exposure is associated with large alpha
changes in either direction rather than consistently larger reductions: of
the seven portfolios with $|\beta_{RMW}| \geq 0.4$, three improve and four
deteriorate.

\paragraph{4(c)}
\begin{enumerate}
\setlength{\itemsep}{2pt}\setlength{\parskip}{0pt}
\item RMW's mean return: The average is 0.242804\% per month, with
$t = 2.966978$. A conventional two-sided $t$-test rejects a zero mean at the
1\% level ($p \approx 0.003102$).
\item Pricing-error improvement: GRS decreases by 10.588612\%, indicating a
modest improvement in the joint criterion. Nevertheless, the augmented-model
$p$-value of $4.213838\times10^{-9}$ and the mixed portfolio-level alpha
changes show that substantial pricing errors remain.
\end{enumerate}
These tests address different questions. The mean-return test asks whether
RMW's average return, its implied premium in the time-series model, differs
from zero; GRS tests whether the 25 intercepts are jointly zero, accounting
for residual covariances and factor moments. A factor with an imprecisely
estimated mean can still reduce pricing errors, because each portfolio's
alpha change depends on its RMW loading and the point estimate of the mean,
not on the precision of that estimate. Conversely, RMW's significant mean
here does not imply that it adequately prices the test assets.

\paragraph{4(d)} The main limitation is overlap with the evidence used to
establish RMW. Fama and French (2015) documented the five-factor model over
July 1963--December 2013 and evaluated it on the same 25 size--B/M
portfolios, so about 80\% of this test's 756 months (606), and the test
assets themselves, overlap the factor's discovery sample; only January
2014--June 2026 (150 months) postdates it. Factor construction adds further
dependence: RMW and the test portfolios are formed from the same U.S.
CRSP/Compustat universe using NYSE-based breakpoints, so they share
underlying stocks. This is therefore not an independent out-of-sample test,
and the historical GRS reduction may overstate performance in new data; this
limits external validity without invalidating the in-sample result.

\paragraph{Prompt 1 --- Step 3(a): Implementation.} ``Write Python code to
estimate the CAPM and augmented RM-RF + RMW time-series regressions for the
25 Size/BE/ME portfolios and run the GRS test. Use the same common sample for
both models, report the GRS statistics and $p$-values, create an alpha
comparison table, and validate the GRS implementation.''

\paragraph{Prompt 2 --- Step 3(b): Interpretation.} ``Using the 25-portfolio
alpha comparison table, identify which portfolios RMW helps price and which
it does not. Focus on changes in absolute alpha, patterns across size and
BE/ME, and whether the results are consistent with my Step 2 hypothesis.''

\subsection{Step 5}

\begin{table}[H]
\centering
\caption{AI interaction Log.}
\footnotesize
\begin{tabular}{p{0.15\textwidth}p{0.2\textwidth}p{0.25\textwidth}p{0.29\textwidth}}
\toprule
Interaction & Prompt/Purpose & What the LLM produced correctly & What required correction or verification \\
\midrule
Step 3(a): Implementation & Prompt 1 above: implement and validate CAPM and
RM-RF + RMW regressions and GRS tests. & The final implementation correctly
handled the two-factor regression and the assignment's GRS framework,
producing GRS statistics, $p$-values, and portfolio alphas. & We verified the
common sample, excess-return construction, $K=1$ versus $K=2$, covariance
denominators, degrees of freedom, and zero-alpha, reorder-invariance, and
scale-invariance checks. No calculation correction was needed; we expanded
the truncated alpha-table display to show all 25 portfolios and added a cell
computing the Sharpe-ratio quantities, RMW mean-test $p$-value, mean RM-RF,
and absolute-alpha totals. \\
\addlinespace
Step 3(b): Alpha interpretation & Prompt 2 above: interpret absolute-alpha
changes, size/BE/ME patterns, and consistency with our hypothesis. & All
quoted figures were accurate, including the largest improvements in
S1\_BM1--S3\_BM1 and deteriorations in S1\_BM2--S1\_BM5. & The description
needed precision: only 9 of 25 portfolios improved, and ``particularly among
smaller firms'' obscured the four largest deteriorations within the
smallest-size quintile. The response omitted the hypothesis comparison and
the role of signed RMW exposures in explaining alpha changes; significant
exposure alone would not establish improvement. \\
\bottomrule
\end{tabular}
\label{tab:cap_log}
\end{table}

\paragraph{Reflection.} Working with the LLM clarified that GRS jointly
evaluates regression intercepts using the residual covariance matrix,
accounting for correlations among pricing errors. Because GRS evaluates the
alphas using the inverse residual covariance matrix, its statistic can fall
even when most individual absolute alphas rise, as in our results. Verifying
covariance denominators and degrees of freedom showed us why invariance
checks alone cannot establish that an implementation is correct. We now
understand GRS as a test of the factors' mean-variance efficiency relative to
the test assets, rather than a count of individually significant alphas.

\end{document}
